\begin{figure}[htbp]
\centering
\begin{tikzpicture}
\begin{axis}[width=0.79\textwidth,height=6cm,
 axis lines=middle,xmin=-1.15,xmax=1.15,ymin=-1.15,ymax=1.3,
 xlabel={$u=x-a$},ylabel={$P$},
 xtick={-1,0,1},ytick={-1,1},
 samples=120,domain=-1.05:1.05,
 legend style={at={(0.03,0.97)},anchor=north west,
 draw=none,fill=white,font=\small},
 label style={font=\small}]
\addplot[arkivblaa,very thick] {x^2};
\addlegendentry{$P=u^2$: lige multiplicitet}
\addplot[arkivrod,very thick,dashed] {x^3};
\addlegendentry{$P=u^3$: ulige multiplicitet}
\end{axis}
\end{tikzpicture}
\caption{Dimensionsløse grundformer ved en rod $x=a$.
En rod med lige multiplicitet berører nul uden fortegnsskift;
en rod med ulige multiplicitet krydser nul. For et generelt polynomium
$(x-a)^mH(x)$ gælder fortegnsargumentet lokalt, hvor $H$ er forskellig
fra nul; $H$ kan ændre både skala og kurvens orientering.}
\label{fig:algebra-multiplicitet}
\end{figure}
